\section{Kesimpulan}
Pada Milestone 1 ini, kelompok kami berhasil mengerjakan perancangan dan implementasi \textit{lexical analyzer} bahasa Arion berbasis DFA, mulai dari pemodelan state, implementasi modul lexer-token-main, hingga pengujian kasus dasar dan \textit{edge cases}. Program yang dihasilkan dapat membaca file input, mengenali token sesuai spesifikasi, serta melaporkan error leksikal dengan konsisten.

Tugas diselesaikan dengan alur yaitu, menerjemahkan spesifikasi ke desain DFA, mengimplementasikan kode secara modular, lalu memvalidasi melalui kumpulan test case. Pelajaran utama yang diperoleh adalah pentingnya dasar teori automata dalam pengembangan program ini, pentingnya pengujian kasus batas/\textit{edge cases}, dan pentingnya komunikasi tim agar integrasi kode serta dokumentasi berjalan lancar.

\section{Saran}

Untuk pengembangan lanjutan, disarankan memperkaya pesan error dengan menambahkan baris/kolom dimana kesalahan terjadi dan menyediakan mode debug transisi DFA. Dari sisi kerja tim, penggunaan alur \textit{branching} yang lebih intensif, \textit{code review} yang lebih terstruktur, dan standar penulisan yang konsisten akan meningkatkan kualitas hasil.

Sebagai langkah berikutnya, proyek dapat dilanjutkan ke analisis sintaksis dan semantik agar lexer yang telah dibangun menjadi fondasi compiler yang lebih lengkap.
